% Beyond the Chomsky Wall: Platonic Representations as the Convergence Point
% of Natural Language and Code
% Formatted for COLM 2026 submission

\documentclass{article}

% COLM style — "submission" for anonymous review
\usepackage[submission]{colm2026_conference}

%% ---- Encoding ----
\usepackage[T1]{fontenc}
\usepackage[utf8]{inputenc}
\usepackage{microtype}
\usepackage{inconsolata}

%% ---- Math ----
\usepackage{amsmath}
\usepackage{amssymb}
\usepackage{amsthm}

%% ---- Tables ----
\usepackage{booktabs}
\usepackage{tabularx}
\usepackage{array}

%% ---- Graphics ----
\usepackage{graphicx}
\usepackage{tikz}
\usetikzlibrary{positioning,arrows.meta}

%% ---- Lists ----
\usepackage{enumitem}

%% ---- Code Listings ----
\usepackage{listings}
\lstset{
  basicstyle=\ttfamily\small,
  breaklines=true,
  frame=single,
  framesep=3pt,
  xleftmargin=5pt,
  xrightmargin=5pt,
  backgroundcolor=\color{gray!5},
  rulecolor=\color{gray!40},
  keywordstyle=\color{blue!70!black},
  commentstyle=\color{green!50!black},
  stringstyle=\color{red!60!black},
  showstringspaces=false,
  tabsize=2
}

%% ---- Custom Theorem Environments ----
\theoremstyle{definition}
\newtheorem{definition}{Definition}[section]
\newtheorem{conjecture}{Conjecture}[section]

%% ---- Convenience Commands ----
\newcommand{\Zsem}{Z_{\text{sem}}}
\newcommand{\Zproc}{Z_{\text{proc}}}
\newcommand{\Zprag}{Z_{\text{prag}}}
\newcommand{\Dtrain}{D_{\text{train}}}

%% ---- Column type for wide tables ----
\newcolumntype{L}{>{\raggedright\arraybackslash}X}

\title{Beyond the Chomsky Wall: Platonic Representations as the\\Convergence Point of Natural Language and Code}

\author{Anonymous}

\date{}

\begin{document}

\maketitle

\begin{abstract}
The Platonic Representation Hypothesis (PRH; Huh et al., ICML 2024) claims that neural networks trained on different modalities converge toward a shared latent structure $Z$. We accept this claim and ask the next question: \emph{does convergence imply communicability?} If two systems---or two cultures---arrive at the same $Z$, can they use it as a shared channel for computational intent? Mathematics suggests yes: the Pythagorean theorem holds universally, yet derivation paths diverge across traditions (axiomatic, constructive, intuitive) and notation persists by convention ($dy/dx$ vs.\ $\dot{y}$). Mathematics functions as a cross-cultural communication medium \emph{despite} this divergence. We ask whether the same can hold for natural language and code. We propose that $Z$ is stratified: $\Zsem$ (what is computed) converges cross-culturally, while $\Zproc$ (how it is derived) and $\Zprag$ (for whom it is framed) remain culturally mediated---and show that existence and communicability are distinct properties. Three waves of natural-language programming (keyword substitution, structured NL, LLM-mediated translation) all operate at the surface level; we reinterpret them as attempts to reach $Z$ that fail because they address projections rather than representations. We formalize disambiguation cost migration via information theory, showing specification complexity is conserved across interfaces, and identify a verification paradox: dissolving the Chomsky wall at input forces its reappearance at verification. A pilot across five languages and 100 operations provides illustrative evidence: cross-lingual probing transfer (P3, 90\% accuracy) and spacing robustness (P7, $R \approx 2.9$--$13.6$) support $\Zsem$ convergence, while cross-lingual NL-NL invariance (P2) fails at the description level even as NL-code alignment succeeds ($R_{\text{code}} > 1$)---demonstrating that representations can converge without descriptions converging, i.e., that existence does not entail communicability. This distinction---between convergence and communicability of shared representations---connects to parallel debates in neuroscience (cross-brain representational similarity), philosophy (the private language problem), and denotational semantics, and we argue it is the central open question for representation-level computing.
\end{abstract}

%% ============================================================
%% 1. INTRODUCTION
%% ============================================================

\section{Introduction}\label{sec:intro}

Mathematics is humanity's most successful Platonic communication medium. The Pythagorean theorem holds in every culture---the \emph{result} is universal. Yet the \emph{derivation} differs: Bourbaki proceeds axiomatically, the Soviet tradition constructively, the Indian tradition intuitively-inductively. The \emph{notation}---Leibniz's $dy/dx$ versus Newton's $\dot{y}$---persists by historical accident. Despite these divergences, mathematicians across cultures reliably communicate through shared ideals. The result is invariant; the path is cultural; the notation is conventional---and yet the system \emph{works} as a communication medium.

But viewing mathematics and natural language as fundamentally separate systems is itself a surface-level distinction. Both are acts of \emph{representation}---ways humans encode their understanding of reality into communicable form. Mathematics constrains its representations to enforce unambiguous denotation; natural language relaxes that constraint to gain expressiveness and efficiency~\citep{piantadosi2012ambiguity}. They are not different kinds of things but different \emph{regimes} of the same representational capacity, trading off precision against flexibility. Code occupies a third regime---formally constrained like mathematics, but operationally directed (it \emph{does} something). From this perspective, convergence across modalities is not surprising; the deeper question is why convergent representations do not automatically yield mutual intelligibility.

Can computation achieve what mathematics has? When a Korean speaker says \emph{``jeongnyeolhara''} and an English speaker says ``sort this list,'' do they access the same computational ideal? If so, can that ideal serve as a shared channel---not merely existing, but being \emph{communicable}?

The Platonic Representation Hypothesis (PRH; Huh et al., ICML 2024)~\citep{huh2024platonic} claims that neural networks trained on different modalities converge toward a shared latent structure $Z$. Ziyin et al.\ (2025)~\citep{ziyin2025proof} prove convergence for deep linear networks. We accept PRH and ask the next question: \textbf{does convergence imply communicability?} This is the computational version of an ancient philosophical problem---whether universals, even if real, are accessible as a shared medium---and it connects to parallel questions in neuroscience, where cross-brain representational similarity~\citep{haxby2011hyperalignment} does not guarantee shared subjective experience.

\paragraph{Position.} This paper does not refute PRH---it \emph{refines} it. We argue that:
\begin{enumerate}[label=(\roman*),nosep]
  \item $Z$ is stratified: $\Zsem$ (what is computed) converges cross-culturally, while $\Zproc$ (derivation path) and $\Zprag$ (communicative frame) remain culturally mediated (\S\ref{sec:stratification}).
  \item Existence and communicability are distinct: $Z$ may converge without being usable as a shared channel (\S\ref{sec:communicability}).
  \item The Chomsky wall between natural language (Type-1) and code (Type-2) constrains surface forms, not representations---but dissolving it at input forces its reappearance at verification (\S\ref{sec:chomsky}, \S\ref{sec:verification}).
  \item Disambiguation cost is conserved: representation-level systems redistribute, rather than eliminate, specification complexity (\S\ref{sec:disambiguation}).
\end{enumerate}

\paragraph{Context: three waves of NL programming.} Attempts to bridge NL and code have come in three waves, all operating at the surface level: \textbf{Wave~1} (keyword substitution, 1960s--) replaces English tokens with native-language equivalents---the grammar remains Type-2. \textbf{Wave~2} (structured NL, 2024--) constrains prose into parseable templates~\citep{core2025,elliott2023sudolang}---new formal grammars that look like English. \textbf{Wave~3} (LLM-mediated translation, 2023--) generates code from free-form NL---the ``vibe coding'' paradigm, but non-deterministic and fundamentally intractable~\citep{icse2026reproducibility}. The QWERTY analogy is instructive: typewriter jams required a specific key layout; the constraint vanished with computers but the layout persisted. Programming syntax may be similar---an engineering constraint maintained by path dependency. But the NL-code gap runs deeper than engineering, into the mathematical structure of grammars (\S\ref{sec:chomsky}). PRH suggests the gap is a property of surface forms, not representations---and the productive question is not whether $Z$ exists, but whether it is communicable.

\paragraph{Contributions.}
\begin{enumerate}[label=(\arabic*),nosep]
  \item \textbf{The communicability problem} as the central question for PRH applied to code: existence of shared representations does not guarantee their utility as a communication medium; mathematics is the precedent, not the guarantee (\S\ref{sec:communicability}).
  \item \textbf{$Z$ stratification} with differential convergence, connecting $\Zsem$ to denotational semantics and identifying $\Dtrain$ dependence as a source of cultural bias even in ``purely computational'' domains (\S\ref{sec:stratification}--\S\ref{sec:dtrain}).
  \item \textbf{Disambiguation bound} formalized via information theory, and a \textbf{verification paradox} showing the Chomsky wall retreats inward rather than dissolving (\S\ref{sec:disambiguation}, \S\ref{sec:verification}).
  \item \textbf{Illustrative pilot evidence} across five languages: P2 (cross-lingual invariance) fails at the description level while NL-code alignment succeeds---directly demonstrating that convergence $\neq$ communicability (\S\ref{sec:pilot}).
  \item \textbf{Cross-disciplinary connections} to neuroscience (representational similarity analysis), philosophy (private language problem), and denotational semantics, situating the communicability gap as a shared challenge across fields (\S\ref{sec:prh}, \S\ref{sec:related}).
\end{enumerate}

\paragraph{Running example.}
Throughout this paper: ``Find the three largest values in the given list.'' This sentence is computationally precise yet linguistically ambiguous---``largest'' lacks type information, comparison criterion, boundary handling, and output ordering. In Type-2 programming, the programmer resolves all ambiguities by writing \texttt{sorted(lst, reverse=True)[:3]}. We track how each wave, each $Z$ layer, and the verification paradox handle these ambiguities.

\paragraph{A cross-lingual grounding example: date parsing.}
Consider the task ``parse this date string into year, month, and day.'' The $\Zsem$ denotation---a function $\text{String} \to (\text{Int}, \text{Int}, \text{Int})$---is culture-invariant. But the surface forms diverge radically: an American writes \texttt{04/03/2026} (MM/DD/YYYY), a Korean writes \texttt{2026.04.03} or \texttt{2026\text{년}4\text{월}3\text{일}}, a German writes \texttt{03.04.2026} (DD.MM.YYYY), and an Arabic speaker may use \texttt{٠٣/٠٤/٢٠٢٦} with Eastern Arabic numerals and right-to-left layout. All denote the same calendar date, yet a parser that hard-codes any single convention fails silently on the others. This is the communicability problem in miniature: the computational meaning converges ($\Zsem$ is identical), the format conventions diverge ($\Zproc$ is culturally mediated), and the pragmatic expectations differ ($\Zprag$---a Korean user expects the year-first output to be the ``natural'' ordering). Wave~1 cannot help (translating keywords does not change date format conventions); Wave~3 introduces non-determinism (an LLM may guess the format or hallucinate). A representation-level system would need to resolve the format ambiguity at $Z$---illustrating disambiguation migration (\S\ref{sec:disambiguation}).

%% ============================================================
%% 2. CURRENT LANDSCAPE
%% ============================================================

\section{Current Landscape: Surface-Form Approaches}\label{sec:landscape}

All three waves operate at the surface-form level. Wave~1 changes terminal symbols but preserves the parse tree. Wave~2 imposes structure on NL, creating new Type-2 grammars---the paradox is that making NL precise enough to execute requires removing the properties that make it natural. Wave~3 works for prototyping but is non-deterministic: our running example produces \texttt{sorted(lst, reverse=True)[:3]}, \texttt{heapq.nlargest(3, lst)}, and a buggy list comprehension across three runs. See Appendix~\ref{app:code} for code examples.

\begin{table*}[t]
\centering
\caption{Dimension comparison across waves and the proposed $Z$-level approach.}
\label{tab:waves}
\small
\begin{tabularx}{\textwidth}{l L L L L}
\toprule
\textbf{Dimension} & \textbf{Wave 1: Keyword} & \textbf{Wave 2: Structured NL} & \textbf{Wave 3: LLM Translation} & \textbf{Proposed: $Z$-Level (this work)} \\
\midrule
Chomsky level & Type-2 (unchanged) & Type-2 (re-imposed) & Type-1 $\to$ Type-2 & Type-1 $\to$ $Z$ $\to$ Type-2 \\
Determinism & Full & Full & Non-deterministic & Deterministic (target) \\
Expressiveness & $=$ base language & $<$ NL & $\approx$ NL & $=$ NL \\
Verification & Standard & Standard & None (stochastic) & At $Z$ (open problem) \\
Disambiguation & By programmer & By programmer & By LLM (non-det.) & By system (at $Z$) \\
\bottomrule
\end{tabularx}
\end{table*}

Both RAG and NL$\to$Code map NL to structured output via learned representations, but code requires exact semantics where RAG tolerates approximation---cosine similarity $0.95$ is fine for document retrieval but a bug for program execution.

%% ============================================================
%% 3. THE PLATONIC REPRESENTATION HYPOTHESIS
%% ============================================================

\section{The Platonic Representation Hypothesis}\label{sec:prh}

\subsection{Overview and Formal Framework}

The naming is not accidental. PRH raises the oldest question in philosophy: do universals exist? Plato's forms were ideal objects that particular instances approximate; PRH claims neural networks discover such forms empirically---that a shared latent structure $Z$ exists, from which images, text, and code are projections. The question ``does $Z$ exist?'' is the computational version of ``do Platonic ideals exist?''---and PRH's answer, supported by Huh et al.\ (ICML 2024)~\citep{huh2024platonic}, is yes, as a statistical convergence. Ziyin et al.\ (2025)~\citep{ziyin2025proof} prove that embedded deep linear networks trained with SGD become ``Perfectly Platonic''---every pair of layers learns the same representation up to rotation. Crucially, they identify six conditions that break convergence (weight decay, label transformation, saddle convergence, input heterogeneity, gradient flow vs.\ SGD, finite-step edge-of-stability), suggesting PRH convergence is the default for well-behaved training but can fail under specific conditions. We introduce semi-formal definitions:

\begin{definition}[Modality Encoder]\label{def:encoder}
An encoder $E_i : M_i \to \mathbb{R}^d$ maps inputs from modality $M_i$ to $d$-dimensional representations; a decoder $D_i : \mathbb{R}^d \to M_i$ maps back.
\end{definition}

\begin{definition}[Representational Convergence]\label{def:convergence}
For semantically equivalent inputs $x \in M_i$ and $y \in M_j$:
\[
\lim_{\text{scale}\to\infty} \mathrm{sim}(E_i(x), E_j(y)) \to 1.
\]
\end{definition}

\begin{definition}[Shared Latent Structure]\label{def:Z}
$Z \approx \lim_{\text{scale}\to\infty} \mathrm{Im}(E_i)$ for all modalities $i$. Critically, $Z = Z(\text{scale}, \Dtrain)$---two models at the same scale trained on different corpora may converge to different $Z$ (\S\ref{sec:dtrain}).
\end{definition}

\begin{definition}[Z Stratification]\label{def:stratification}
$Z$ decomposes into: $\Zsem$ (what is computed, cross-culturally convergent), $\Zproc$ (how it is derived, culturally mediated), $\Zprag$ (for whom it is framed, culturally specific).
\end{definition}

\begin{definition}[The Chomsky Wall]\label{def:wall}
The computational boundary between Type-2 and Type-1. Surface-level if it constrains only parsing, not representations.
\end{definition}

\paragraph{Central Claim.}
We do not contest PRH---we extend it. PRH establishes that $Z$ converges; we ask whether convergent $Z$ is \emph{communicable}. If PRH holds for code as a modality, the Chomsky wall constrains only the projection $D_i : Z \to M_i$, not $Z$ itself. Operations at $Z$ are not subject to Chomsky hierarchy constraints because $Z$ is continuous---it has no grammar to classify. But the absence of a grammar wall at $Z$ does not guarantee that two agents accessing $Z$ through different encoders (languages, cultures, training data) arrive at the same point---or that they can verify they have.

\subsection{Connection to Denotational Semantics}

Denotational semantics (Scott \& Strachey, 1971)~\citep{scott1971semantics} maps programs to mathematical objects in a semantic domain $\mathcal{D}$, independent of syntax. We observe a structural parallel: $\Zsem$ is a neural approximation of $\mathcal{D}$.

\begin{conjecture}\label{conj:denotational}
$\lim_{\text{scale}\to\infty} \Zsem(E(d)) \approx [\![ P ]\!]$ (up to continuous bijection), where $d$ is a NL description and $[\![ P ]\!]$ is the denotation of program $P$.
\end{conjecture}

This explains why $\Zsem$ converges cross-culturally (denotational semantics is culture-invariant), why convergence is stronger for computational domains (well-defined denotations), and why compositionality is the hard problem (denotational compositionality requires homomorphic structure NL does not guarantee). For our running example: the denotation of \texttt{sorted(lst, reverse=True)[:3]} is a function $\text{List}[\text{Comparable}] \to \text{List}[\text{Comparable}]$; PRH claims ``find the three largest values'' converges to the same point in $\Zsem$.

\subsection{Evidence for Code as a Modality}\label{sec:evidence}

\begin{itemize}[nosep]
  \item \textbf{Aligned embeddings.} UniXcoder achieves 77.1\% MRR on NL-code search across six languages.
  \item \textbf{Cross-model transferability.} Steering vectors from 7B models control 70B model behavior (ACL 2025)~\citep{chen2025crossmodel}---$Z$ structure is recoverable across architectures.
  \item \textbf{Language-independent code semantics.} LLMs develop language-independent semantic representations of code, with syntax-specific components in early layers feeding shared semantic components in middle layers (ICSE 2025 NIER)~\citep{beyondsyntax2025}.
  \item \textbf{The Naturalness Hypothesis.} Software is statistically natural---code cross-entropy approaches that of English text (Hindle et al., 2012)~\citep{hindle2012naturalness}, a precondition for PRH convergence.
  \item \textbf{Translation emergence.} LLMs trained primarily on NL generate code without code-specific training; cross-language synergies follow scaling laws (Cassano et al., 2025)~\citep{cassano2025scaling}.
  \item \textbf{Bidirectional transfer.} Training on NL documentation improves code completion, and vice versa.
\end{itemize}

We note that PRH for code remains a conjecture---the evidence above is consistent with convergence but does not prove it. Our contribution is not to establish PRH for code but to derive its consequences and provide initial empirical tests (\S\ref{sec:pilot}).

\subsection{The Stratification of Z}\label{sec:stratification}

Mathematics already embodies the structure PRH predicts for $Z$. The Pythagorean theorem holds universally---but a Bourbaki algebraist proves it axiomatically, a Soviet constructivist computes it explicitly, and an Indian mathematician reaches it through geometric intuition. The result is identical; the derivation is culturally mediated; the notation ($a^2 + b^2 = c^2$ vs.\ equivalent formulations) is a surface convention. If PRH's $Z$ is real, it should exhibit this same stratification: universal at the level of what is computed, culturally distributed at the level of how, and conventional at the level of how it is expressed. The fact that mathematics functions as a cross-cultural communication medium---despite culturally divergent derivation paths---is direct evidence that Platonic ideals, when they exist, can serve as a shared channel. The question for NL-code convergence is whether $\Zsem$ achieves this same communicability for computational intent. Similarly, \texttt{sort(array)} produces the same output regardless of functional or imperative derivation.

\begin{table}[t]
\centering
\caption{$Z$ stratification: layers, content, and convergence.}
\label{tab:stratification}
\small
\begin{tabular}{lll}
\toprule
\textbf{Layer} & \textbf{Content} & \textbf{Convergence} \\
\midrule
$\Zsem$ & What is computed & Cross-culturally convergent \\
$\Zproc$ & How it is derived & Culturally mediated \\
$\Zprag$ & For whom framed & Culturally specific \\
\bottomrule
\end{tabular}
\end{table}

\paragraph{Running example.}
$\Zsem$: select top-3 by magnitude (culture-invariant). $\Zproc$: Korean developer may filter (``do not exclude large values''), Haskell developer may write \texttt{take 3 . sortBy (flip compare)}. $\Zprag$: the suffix \emph{-ara} (a Korean imperative) is a Korean imperative frame.

\subsection{Multilingual Convergence and $\Dtrain$ Dependence}\label{sec:dtrain}

PRH implies a single convergence point. Multilingual LLMs complicate this. Two models compete:

\paragraph{Model A: The multilingual human.}
The LLM code-switches between culturally-mediated representations---Korean input activates a Korean-inflected $Z$; English input activates an Anglo-American one.

\paragraph{Model B: The Platonic oracle.}
The LLM contains a culture-invariant $Z$; cultural variations are decoding artifacts.

Evidence increasingly favors Model~A for judgment, Model~B for structure: 9--56pp moral judgment gaps across languages including preference reversals~\citep{vida2024moral,agarwal2024ethical}, GPT-4o displays Cultural Frame Switching with distinct ``personalities'' per language (COLING 2025)~\citep{ahn2025coling}, and reasoning-language effects contribute twice the variance of input-language effects~\citep{li2026reasoning}. But shared grammatical features persist across typologically diverse languages~\citep{brinkmann2025shared}, a semantic hub appears in middle layers~\citep{wu2025semantichub}, and language identity is encoded via sparse dimensions separate from semantic content~\citep{zhong2025sparse}.

\paragraph{Resolution.}
$\Zsem$ converges (Model~B); $\Zproc$ and $\Zprag$ remain culturally distributed (Model~A). The Aristotelian critique (Gr\"{o}ger et al., 2026)~\citep{groger2026aristotelian} refines further: what converges is local neighborhood topology (preserved neighborhoods), not global geometry (preserved distances). $Z$ may be a topological rather than geometric structure.

\paragraph{$\Dtrain$ dependence.}
$Z$ is not determined by scale alone---it depends on training data cultural composition: $Z = Z(\text{scale}, \Dtrain)$. Even ``purely computational'' domains carry implicit cultural defaults: sort order conventions, null handling semantics, error semantics (Python exceptions vs.\ Go explicit errors vs.\ C undefined behavior), and number formatting. The clean $\Zsem$/$\Zproc$ separation is an idealization:
\[
\Zsem(\text{observed}) = \Zsem(\text{ideal}) + \mathrm{bias}(\Dtrain).
\]
The sovereign AI movement (Korea's five-consortium initiative, Japan's ABCI 3.0, SEA-LION) implicitly acknowledges this---if $Z$ were truly Platonic, culturally-native models would be unnecessary. The market is voting for Model~A. See Appendix~\ref{app:dtrain} for detailed analysis.

\subsection{Why Does Z Converge?}\label{sec:whydoes}

PRH attributes convergence to a shared statistical model of \emph{reality}~\citep{huh2024platonic}. We propose an alternative: convergence reflects the cognitive homogeneity of its sole training species.

\paragraph{Species-level constraint.}
Homo sapiens is among the most genetically homogeneous large mammals. A single group of 50 chimpanzees harbors more nucleotide diversity than all 8 billion humans~\citep{pradomartinez2013apes}. A severe bottleneck ${\sim}930{,}000$ years ago reduced the breeding population to ${\sim}1{,}280$ individuals, erasing an estimated 65.85\% of genetic diversity~\citep{hu2023bottleneck}. All LLM training data is produced by this genetically constrained species; the apparent universality of $Z$ may reflect the narrow bandwidth of human perception rather than a Platonic structure of the world.

\paragraph{WEIRD double constraint.}
Internet text further narrows this sample. Henrich et al.\ (2010)~\citep{henrich2010weird} showed that WEIRD (Western, Educated, Industrialized, Rich, Democratic) subjects---from whom the majority of training data originates---are ``among the least representative populations one could find for generalizing about humans.'' $Z$ is thus doubly constrained: by species-level cognitive architecture and by the cultural subsample that dominates $\Dtrain$.

\paragraph{Biological precedent.}
Biology provides abundant cases where signals converge \emph{without} communicability---precisely the pattern we argue holds for $Z$:
\begin{itemize}[nosep]
  \item \textbf{Quorum sensing.} Autoinducer-2 (AI-2) is produced by $>$50\% of sequenced bacterial species, yet each species uses different receptors and triggers different gene cascades~\citep{waters2005quorum}. Same representation, species-specific interpretation.
  \item \textbf{Pheromone reuse.} The same sesquiterpene compound functions as alarm pheromone in one species and sex pheromone in another. Identical signal, opposite semantics.
  \item \textbf{Heterospecific eavesdropping.} Fairy-wrens decode other species' alarm calls only in sympatry---where they co-occur and co-train~\citep{magrath2015eavesdropping}. Communicability requires alignment beyond signal convergence.
\end{itemize}
These are not edge cases; convergence without communicability is the biological norm. Vervet monkeys' alarm calls~\citep{seyfarth1980vervet}---acoustically distinct signals for leopards, eagles, and snakes---show that even within a species, signal-meaning mappings must be \emph{learned}, not inherited. The convergent signal is necessary but not sufficient.

\paragraph{Implication for PRH.}
If convergence reflects $Z(\text{Homo sapiens})$ rather than $Z(\text{reality})$, the communicability gap is deeper than PRH suggests. It is not merely that convergent representations are hard to \emph{access} across cultures (a practical problem); it is that what converges may be a species-specific projection of reality, further distorted by a WEIRD subsample (a structural problem). Model collapse from recursive training~\citep{shumailov2024collapse} compounds this: tail distributions---including minority cultural data---vanish first.

\subsection{The Communicability Problem}\label{sec:communicability}

The preceding sections establish that $Z$ may exist (PRH convergence), that it stratifies ($\Zsem$ converges while $\Zproc$ does not), and that this stratification depends on training data. But existence is not sufficient for utility. \textbf{The central question of this paper is whether $Z$ can serve as a communication medium}---whether a Korean speaker's computational intent, encoded to $Z$, is decoded by the system as the same intent.

This question has deep roots. In neuroscience, Hyperalignment~\citep{haxby2011hyperalignment} demonstrates that different brains produce geometrically similar representations when processing identical stimuli---a biological analogue of PRH convergence. Representational Similarity Analysis~\citep{kriegeskorte2008rsa} quantifies this convergence using the same mathematical tools (second-order similarity matrices) as CKA and probing in AI. Yet shared representational geometry does not resolve Wittgenstein's private language problem~\citep{wittgenstein1953}: even if your neural representation of ``red'' is geometrically similar to mine, we cannot verify that we \emph{experience} the same red. ``Is my $\Zsem$ your $\Zsem$?'' is the computational version of this question.

Mathematics provides the strongest evidence that communicability \emph{is} achievable despite cultural divergence: $a^2 + b^2 = c^2$ communicates identically across cultures because the denotation is unambiguous. But even mathematical notation is not fully culture-free: the number 4 carries death associations in CJK cultures ($\text{四} \approx \text{死}$), 13 is avoided in Western contexts, and 8 connotes prosperity in Chinese ($\text{八} \approx \text{發}$). A \texttt{sort(array)[:4]} is denotationally identical everywhere, but its connotation leaks $\Zprag$---the ``purely computational'' carries cultural residue. Mathematics achieves communicability precisely by \emph{bypassing} natural language---encoding denotations directly in formal notation---yet even this bypass is imperfect. Natural language does not even attempt the bypass: ``find the largest values'' carries cultural defaults invisible to the speaker.

A structural asymmetry deepens this gap. Code is \emph{modality-complete}: its text encodes the entirety of its meaning, with no reliance on intonation, gesture, or shared background. Natural language is \emph{modality-incomplete}: meaning depends on non-verbal channels that embedding models cannot access. High-context languages~\citep{hall1976beyond} amplify this incompleteness: Korean routinely omits subjects recoverable from social context, and mimetic expressions like \emph{salgeumsalgeum} (``stealthily'') encode embodied sensory meaning with no clean English mapping. Embedding models trained on text alone capture only the explicit channel---for code this is everything; for NL, especially high-context NL, it is a lossy projection. Recent work on World English varieties confirms this asymmetry empirically: AI-generated text maintains high semantic similarity while systematically erasing pragmatic and cultural markers, with pragmatic markers $1.9\times$ more vulnerable than lexical ones~\citep{culturalghosting2026}.

Our pilot experiment (\S\ref{sec:pilot}) sharpens this point: \emph{description-level} cross-lingual invariance does not follow the pattern PRH predicts for \emph{execution-level} convergence. NL descriptions of computational operations are \emph{less} invariant across languages than descriptions of judgment operations---the opposite of what $\Zsem$ convergence would suggest. This is direct evidence for the communicability gap: $\Zsem$ may converge at the level of what programs \emph{do}, even when NL descriptions of what they do diverge. The gap is located not in $Z$ itself but in the NL-to-$Z$ encoding path---precisely where cultural mediation occurs, and where modality incompleteness is most acute.

\subsection{The Convergence-Determinism Gap}

PRH gives $\mathrm{sim} \to 1 - \delta$, but execution requires $\delta = 0$. Three resolutions:
\begin{enumerate}[nosep]
  \item \textbf{Discretization at $Z$}---equivalence classes, formally equivalent to constructing a new grammar over $Z$.
  \item \textbf{Verification after projection}---accept $\delta > 0$ at $Z$, demand $\delta = 0$ at output.
  \item \textbf{Probabilistic execution}---statistical guarantees (probability $1 - 10^{-9}$).
\end{enumerate}
This gap is the central technical obstacle.

\subsection{Falsifiable Predictions}\label{sec:predictions}

\begin{description}[nosep,font=\bfseries]
  \item[P1 (Scale-Convergence).] NL-code cosine distance decreases monotonically with model scale.
  \item[P2 (Cross-Lingual Semantic Invariance).] \emph{jeongnyeolhara} (Korean: sort) and ``sort this'' (English) are closer in $Z$ than ``sort this'' and ``reverse this'' in English. This should not hold for judgment-involving operations. \emph{Extension:} this invariance should degrade along a dialect continuum---cross-dialectal $R$ (e.g., standard Korean vs.\ Jeju) should fall between within-language $R$ and cross-lingual $R$, revealing that the communicability gap is continuous, not categorical.
  \item[P3 (Stratification Separability).] Probing classifiers for \emph{what is computed} generalize cross-lingually; probes for \emph{how} show language-specific patterns.
  \item[P4 (Domain-Dependent Determinism).] Constrained decoding achieves higher consistency for computational than judgment tasks.
  \item[P5 (Disambiguation Migration).] Total disambiguation effort is conserved---syntax errors migrate to semantic clarification queries.
  \item[P6 ($\Dtrain$ Sensitivity).] $\Zsem$ divergence across models correlates with $\Dtrain$ cultural divergence, even for ``purely computational'' operations with implicit criteria.
  \item[P7 (Spacing Robustness).] $Z$ distance between spacing variants of the same sentence $<$ $Z$ distance between semantically different sentences with identical spacing (see Appendix~\ref{app:tokenization}).
  \item[P8 (Keyboard Transliteration Fragility).] Typing Korean on an English-layout keyboard without switching (e.g., \emph{dkssudgktpdy} $=$ 안녕하세요) produces a string with identical intent but radically different $Z$. Search engines resolve this via layout mapping; embedding models do not. $Z$ distance between transliterated and correct input should be near-maximal---the most destructive surface variant, and a direct test of whether $Z$ encodes intent or surface form.
\end{description}

%% ============================================================
%% 4. REINTERPRETING THE CHOMSKY WALL
%% ============================================================

\section{Reinterpreting the Chomsky Wall}\label{sec:chomsky}

\subsection{The Wall Is Real---at the Surface}

\begin{table}[t]
\centering
\caption{Chomsky hierarchy: grammar types, parsing complexity, and examples.}
\label{tab:chomsky}
\small
\begin{tabular}{llll}
\toprule
\textbf{Type} & \textbf{Grammar} & \textbf{Parsing Complexity} & \textbf{Examples} \\
\midrule
3 & Regular & $O(n)$ & Regex \\
2 & Context-free & $O(n^3)$, deterministic & Python, C, Java \\
1 & Context-sensitive & PSPACE-complete & Natural language \\
0 & Unrestricted & Undecidable & Turing machines \\
\bottomrule
\end{tabular}
\end{table}

No deterministic algorithm resolves Type-1 ambiguity in polynomial time. Type-2 constraints also function as a disambiguation mechanism: every valid expression has exactly one parse.

\subsection{The Wall Dissolves---at the Representation Level}

LLMs map input sequences to high-dimensional representations via attention over full context. ``if x is greater than 5'' (Korean NL), \texttt{if x > 5:} (Python), \texttt{(> x 5)} (Lisp) all map to nearby points (Semantic Hub Hypothesis). Language identity is encoded via sparse dimensions orthogonal to semantic content~\citep{zhong2025sparse}. The Chomsky hierarchy describes parsing complexity of surface forms---it says nothing about representing meaning.

\paragraph{Surface-form boundaries.}
Between NL and $Z$ exists a spectrum of representational directness: NL (linguistic, culturally mediated) $\to$ emoji/pictographs (iconic, semi-universal) $\to$ mathematical symbols (formal, universal) $\to$ code operators $\to$ $Z$ (continuous). Emoji are natural experiments in $\Zsem$---near-universal yet exhibiting cultural variation that mirrors $Z$ stratification (e.g., the prayer/gratitude/high-five emoji). See Appendix~\ref{app:emoji}. Tokenization introduces a bottleneck: spacing conventions (Korean \emph{ttieo-sseugi}, Chinese segmentation) leak through tokenizers into $Z$, and a ``token tax'' (Petrov et al., 2024)~\citep{petrov2024tokenizers} means non-Latin languages require 2--15$\times$ more tokens for the same content. See Appendix~\ref{app:tokenization}.

\subsection{Disambiguation Migration}\label{sec:disambiguation}

Removing the Chomsky wall does not remove the need for disambiguation---it migrates the cost. In Type-2 programming, the programmer disambiguates at write time via grammar constraints; in representation-level execution, the system disambiguates at $Z$.

We conjecture a \emph{disambiguation bound}: any system accepting ambiguous input must resolve ambiguity somewhere, and the total specification cost is bounded below across architectures. We use ``bound'' rather than ``conservation'' deliberately: the inequality in Eq.~\ref{eq:conservation} establishes a lower bound on specification complexity, not an equality. A true conservation law would require proving that no interface can reduce total disambiguation cost---we claim only that it cannot be eliminated, while acknowledging that different interfaces may achieve different total costs within the bound. This is a design heuristic, not a formal theorem---but it suggests that representation-level systems should not promise to eliminate programming difficulty, only to relocate it from syntactic precision to semantic verification.

The practical question is whether the relocation is net beneficial: is it easier for a human to learn \texttt{if x > 5:} or to verify that the system correctly interpreted ``exclude large values''? For professional developers, the current paradigm may be more efficient. For the billions who can express computational intent in natural language but cannot program, the tradeoff favors representation-level systems even if disambiguation cost is conserved.

\paragraph{Running example.}
``find the three largest values'' contains at least four ambiguities: (1)~type---what kind of values? (2)~comparison criterion---``the largest'' by what measure? (3)~boundary---fewer than three elements? (4)~output ordering---sorted result? In Type-2, \texttt{sorted(lst, reverse=True)[:3]} resolves all four silently---the type system handles~(1), comparison operator handles~(2), slicing handles~(3), and sort determines~(4). At $Z$, each must be resolved by inference, query, or default.

\paragraph{Information-theoretic formalization.}
For task $T$ with unique behavior $B(T)$, the specification complexity $K(T)$ is invariant across interfaces:
\begin{equation}\label{eq:conservation}
K(T) \leq H(S) + H(T|S) + O(1)
\end{equation}
where $H(S)$ is surface-form entropy and $H(T|S)$ is residual ambiguity. Type-2 interfaces maximize $H(S)$ (syntactic precision) and minimize $H(T|S)$. NL interfaces minimize $H(S)$ but face $H(T|S) > 0$---remaining bits must be provided through disambiguation. This connects to Piantadosi et al.\ (2012)~\citep{piantadosi2012ambiguity}, who argue ambiguity is a feature of efficient NL communication---but resolution cost cannot be avoided when exact computation is required. Via Rate-Distortion Theory~\citep{shannon1959coding}: there is a minimum specification rate below which semantic error exceeds acceptable thresholds. Type-2 operates at zero distortion with high rate; NL operates at lower rate with nonzero distortion corrected downstream.

\paragraph{Implication.}
Representation-level systems do not reduce total specification information---they redistribute it: less from the user upfront (lower syntactic cognitive load), more from the system at $Z$ (higher inference/verification cost).

\subsection{Why Translation Fails but Representation Might Work}

Wave~3 accepts NL (Type-1), translates to code (Type-2), executes. The translation step is non-deterministic and does not preserve semantics~\citep{cheung2025translation}. A representation-level approach instead: (1)~accept NL, (2)~map to $Z$, (3)~execute from $Z$ directly or project to code deterministically via constrained decoding.

%% ============================================================
%% 5. TOWARD REPRESENTATION-LEVEL EXECUTION
%% ============================================================

\section{Toward Representation-Level Execution}\label{sec:toward}

A representation-level system would: (1)~accept any surface form, (2)~normalize (spacing, OCR, ASR), (3)~encode to $Z$, (4)~verify and disambiguate at $Z$, (5)~execute at $Z$ or project to Type-2. Several pieces exist: constrained decoding (Outlines, LMQL) makes $Z \to$ Type-2 deterministic; Large Concept Models~\citep{meta2024lcm} operate on SONAR sentence embeddings; LPN~\citep{macfarlane2025lpn} and COCONUT~\citep{hao2024coconut} demonstrate latent execution; neuro-symbolic systems (Scallop~\citep{li2023scallop}, Lobster~\citep{huang2026lobster}) bridge continuous representations and discrete logic. The gap is unifying these with PRH-predicted convergence.

Key open problems include: determinism at $Z$ (convergence $\neq$ identity), verification without Type-2 projection, compositionality (NL composes loosely; code composes strictly), and cultural invariance (when $\Zsem$ carries cultural inflection, the ideal is a family $Z(c)$, not a single $Z$).

\subsection{Illustrative Pilot}\label{sec:pilot}

We test four predictions using 100 operations (50 computational + 50 judgment) described in five languages (English, Korean, Chinese, Arabic, Spanish) with dialectal variants (e.g., standard Korean vs.\ Gyeongsang/Jeolla/Jeju; MSA vs.\ Egyptian/Gulf/Levantine Arabic), embedded through multilingual sentence-transformers (MiniLM-L12, E5-large) and UniXcoder. The discriminability ratio $R = d_{\text{inter}}/d_{\text{intra}}$ tests whether cross-lingual same-operation similarity exceeds within-language different-operation similarity. Full protocol in Appendix~\ref{app:pilot}.

\paragraph{P1 (Scale-Convergence): Architecture-dependent.}
When testing across different model families (MiniLM 118M, mpnet-base 278M, E5-large 560M), scale shows no effect: Spearman $\rho = -0.500$, $p = 0.67$. However, within the E5 architecture family---same architecture, three scales---$R_{\text{total}}$ increases monotonically: E5-small (384d) $R = 2.25$, E5-base (768d) $R = 2.67$, E5-large (1024d) $R = 2.81$ ($\rho = 1.0$, $n = 3$). At the same dimensionality (1024d), BGE-M3 achieves $R_{\text{total}} = 5.58$ versus E5-large's $2.81$---a $2\times$ difference attributable to architecture and training data alone. \textbf{P1 is supported within architecture family but dominated by architecture choice across families.} This suggests that convergence is real but its rate depends more on model design than on raw scale---a refinement of PRH's scale prediction.

\paragraph{P2 (Cross-Lingual Invariance): Not supported---but informative.}
Judgment operations show \emph{higher} cross-lingual invariance than computational operations ($R_J = 3.79$ vs.\ $R_C = 1.91$ for MiniLM; $p=1.0$), consistently across all five models tested---the opposite of our prediction. Computational operations have higher $d_{\text{intra}}$ (0.295 vs.\ 0.191, $p < 0.0001$), meaning they are \emph{less} invariant at the description level. This is robust: no model showed $R_C > R_J$.

\paragraph{What P2 reveals about $Z$ stratification.}
The P2 prediction was wrong. We present two competing interpretations of the $R_J > R_C$ finding, which we regard as \emph{equally viable} given the current evidence.

\paragraph{Interpretation A: Convergence is real but the test cannot reach it.}
Three complementary hypotheses support this reading. \emph{(i) English-pivot alignment}: Multilingual transformers process non-English input through intermediate English-like representations~\citep{wendler2024llamas}. Judgment operations use abstract vocabulary (``evaluate,'' ``prioritize'') whose English glosses are near-identical across languages, yielding high $R_J$; computational operations use domain-specific terminology that maps to diverse English surface forms, lowering $R_C$. \emph{(ii) Description vs.\ execution semantics}: Following SemCoder's~\citep{wan2024semcoder} three-level distinction (description, execution, behavior), our embedding-based $R$ operates at the description level---surface glosses may diverge even when execution semantics converge. \emph{(iii) $\Dtrain$ mediation}: English-dominant training causes computational operations, with more domain-specific vocabulary, to accumulate larger $\Dtrain$-dependent bias~\citep{wu2025semantichub}. Under Interpretation~A, P2 tests the NL-to-$Z$ encoding interface rather than $Z$ itself.

\paragraph{Interpretation B: Computational semantics may not converge uniformly.}
An equally plausible reading is that $\Zsem$ itself does not converge for computational operations. Computational semantics carry genuinely culture-dependent conventions: collation orders (ascending = alphabetical in English, stroke-count in CJK, JIS code in Japanese), null handling (SQL \texttt{NULL} $\neq$ Python \texttt{None} $\neq$ Go explicit error), date format conventions (\S\ref{sec:intro}), and directionality assumptions. Judgment operations, by contrast, draw on vocabulary already cross-lingually calibrated through international business, diplomacy, and academic discourse---their higher $R_J$ may reflect genuine cross-cultural semantic alignment rather than measurement artifact. If computational $\Zsem$ does not converge uniformly, then PRH's single-convergence-point prediction must be qualified by operation type, and the $\Zsem$/$\Zproc$/$\Zprag$ stratification may not carve $Z$ at its joints---the three-layer model may be an over-specification, or the layers may interact in ways that prevent clean empirical separation.

We cannot distinguish Interpretation~A from Interpretation~B with the current pilot; doing so requires execution-level probing (embedding code behavior rather than NL descriptions). The $R_J > R_C$ asymmetry is, however, a result in its own right---one that constrains future Z-level system design regardless of which interpretation proves correct.

\emph{Note on P3 consistency}: The English-pivot hypothesis (Interpretation~A.i) raises a methodological concern that applies symmetrically to P3. If multilingual transformers route through English-like representations, then P3's high cross-lingual probe transfer (90\% category, 86\% operation) may partly reflect this architectural property rather than genuine $\Zsem$ generalization. We flag this as an open confound that future work must address with monolingual model controls.

\paragraph{P7 (Spacing Robustness): Strongly supported.}
Korean spacing variants cluster ${\sim}3\times$ closer than semantically different operations ($R_{\text{spacing}} = 2.90$, $p < 0.001$). Punctuation robustness is even stronger: $R_{\text{punct}} = 13.6$. The one outlier is UPPERCASE (drift = 0.192), acting as a pragmatic signal---evidence that $\Zprag$ is encoded in surface cues even when $\Zsem$ is unchanged.

\paragraph{NL-Code Alignment: Supports PRH.}
Embedding NL descriptions alongside Python code through UniXcoder yields $R_{\text{code}} = 1.043$ (95\% CI: $[1.021, 1.066]$, $p < 0.0001$ by permutation test, $n = 10{,}000$; Cohen's $d = 0.33$). MiniLM shows a larger effect: $R_{\text{code}} = 1.182$ (95\% CI: $[1.147, 1.219]$, $p < 0.0001$, $d = 0.94$). NL maps closer to its corresponding code than to other code, though the UniXcoder effect is small in absolute terms. This resolves the P2 ambiguity---NL descriptions diverge across languages (P2 fails), yet each maps closer to its code equivalent (NL-code alignment holds). \textbf{Description-level invariance and execution-level convergence are distinct phenomena}---the central empirical finding.

\paragraph{P3 (Stratification Separability): Supported, with caveats.}
Linear probes trained on English embeddings achieve 90\% mean accuracy on category transfer and 85.8\% on 100-way operation ID across non-English languages (chance: 50\% and 1\%). Korean transfers worst (80\%/66\%), consistent with $\Dtrain$ effects. However, these results must be interpreted alongside the English-pivot concern raised in Interpretation~A.i: if multilingual transformers route non-English inputs through English-like intermediate representations~\citep{wendler2024llamas}, then high cross-lingual probe transfer may partly reflect this architectural property rather than genuine $\Zsem$ generalization. Distinguishing these accounts requires testing with monolingual models that lack English-pivot architecture. We therefore report P3 as \emph{consistent with} $\Zsem$ generalization but not conclusive evidence for it.

%% ============================================================
%% 6. DISCUSSION
%% ============================================================

\section{Discussion}

\subsection{The Verification Paradox}\label{sec:verification}

If verification requires formal methods, and formal methods require Type-2 input, the system must project from $Z$ to Type-2 for verification---reintroducing the Chomsky wall at the verification stage. The wall is dissolved at input but reassembled at verification. Near-term systems will use a hybrid: $Z$ provides expressiveness; Type-2 projection provides verifiability. The paradox is not fatal but constrains the deployment path: the wall retreats inward, from the user-facing boundary to the system-internal verification layer.

\paragraph{The AI-AI communication corollary.}
A revealing special case: when agents communicate \emph{with each other} rather than with humans, the NL-to-$Z$ encoding path is eliminated entirely. Agent~A's output is already in $Z$; Agent~B consumes $Z$ directly. If the communicability gap resides in the NL-to-$Z$ interface rather than in $Z$ itself, then $Z$-to-$Z$ communication should exhibit no gap---convergence \emph{would} imply communicability within the latent space. Precedents include Facebook's negotiation bots~\citep{lewis2017negotiation}, which spontaneously developed abbreviated language incomprehensible to humans but efficient between agents, and COCONUT~\citep{hao2024coconut}, which performs reasoning in continuous latent space without decoding to NL tokens. This yields a three-way trade-off:

\begin{center}
\footnotesize
\begin{tabular}{@{}lccc@{}}
\toprule
\textbf{Channel} & \textbf{Efficiency} & \textbf{Auditability} & \textbf{Gap} \\
\midrule
Human--AI (NL) & Low & High & Present \\
AI--AI ($Z$) & High & None & Absent \\
AI--AI $\to$ Human summary & Medium & Partial & \textbf{Re-emerges} \\
\bottomrule
\end{tabular}
\end{center}

The third row is critical: even if agents communicate losslessly in $Z$, the moment the result is projected to NL for human consumption, the communicability gap re-emerges. The gap does not disappear; it \emph{migrates} from input to output. This is structurally identical to the Verification Paradox: the wall retreats inward. For multi-agent systems, this implies a fundamental design tension between communication efficiency and human oversight---a tension that spatial-metaphor interfaces (showing \emph{patterns} rather than \emph{content}) may partially resolve by providing ambient awareness without requiring full NL decoding.

Notably, agent-to-agent NL communication represents a \emph{regression} in computing history. Machine-to-machine communication progressed from binary protocols (1950s--90s) through structured embeddings (Word2Vec, GloVe: 300-dimensional vectors encoding semantics more efficiently than NL) to TCP/IP packet switching---none of which required human-readable intermediaries. Current LLM agents perform an unnecessary round-trip: $Z \to \text{NL} \to Z$, decoding internal representations to text only to have the receiving agent re-encode them. This is analogous to requiring all network packets to be formatted as human-readable letters. The architectural resolution is the same one networking adopted: agents communicate in structured $Z$-derived protocols (the ``packet''), and human-facing interfaces render patterns (the ``browser'') without requiring full NL decoding of every exchange.

\subsection{Cultural Boundary}

$\Zsem$ is culture-invariant for purely computational domains; $Z$ carries cultural inflection for judgment-involving domains---requiring $Z(c)$ rather than $Z$. This inflection is amplified in high-context cultures~\citep{hall1976beyond}, where meaning relies on non-verbal channels---subject omission, social register, embodied reference---that text-only models cannot capture~\citep{park2024pragmatic}. The modality incompleteness of NL thus varies across cultures, making the communicability gap a function of both $\Dtrain$ and the source culture's context dependence. Furthermore, the gap is not categorical but \emph{continuous}: dialectal variants within a language occupy intermediate positions, as shown by dialect-aware embedding benchmarks~\citep{faisal2025swan} and evidence that LLMs treat dialects as stigmatized deviations rather than points on a continuum~\citep{bui2025dialect}. $\Dtrain$ dependence creates ethical risks: computational colonialism (English-derived norms as ``neutral'' $Z$), consent opacity, and cultural feedback loops ($Z$-projected outputs re-enter $\Dtrain$). Korean programmers face a double translation: thought (Korean) $\to$ code (English syntax) $\to$ execution; representation-level systems could eliminate both steps~\citep{li2026reasoning}. The biological precedents in \S\ref{sec:whydoes} suggest this boundary is not unique to AI: convergent signaling without communicability is the norm across biological systems, from bacterial quorum sensing to primate alarm calls. The NL-code communicability gap may be an instance of a more general principle: convergent representations arise readily under shared environmental pressure, but communicability requires additional alignment that convergence alone does not provide. See Appendix~\ref{app:ethics} for full analysis.

% Limitations in standalone section after Conclusion

%% ============================================================
%% 7. RELATED WORK
%% ============================================================

\section{Related Work}\label{sec:related}

\paragraph{PRH and convergence.}
Huh et al.\ (ICML 2024)~\citep{huh2024platonic} established convergence; Ziyin et al.\ (2025)~\citep{ziyin2025proof} proved it for deep linear networks; Gr\"{o}ger et al.\ (2026)~\citep{groger2026aristotelian} showed global convergence is confounded but local topology persists. We contribute a biological/anthropological critique (\S\ref{sec:whydoes}): convergence may reflect species-level cognitive homogeneity---Homo sapiens' unusually low genetic diversity~\citep{pradomartinez2013apes,hu2023bottleneck}---rather than a Platonic structure of reality. Biology provides ample precedent for convergent signals without communicability~\citep{waters2005quorum,seyfarth1980vervet,magrath2015eavesdropping}.

\paragraph{Multilingual representations.}
The Semantic Hub Hypothesis~\citep{wu2025semantichub} demonstrates shared representations; middle-layer alignment is strongest across 1,000+ language pairs~\citep{middlelayer2025}. However, multilingual $\neq$ multicultural (Rystr{\o}m et al., 2025)~\citep{rystrom2025multicultural}: LLMs produce different moral judgments across languages, reflect creator ideology~\citep{buyl2025ideology}, and Sapir-Whorf effects are measurable in LLM outputs across 13 languages~\citep{ray2025sapirwhorf}.

\paragraph{NL programming and vibe coding.}
Wave~3's non-determinism is fundamentally intractable (ICSE 2026)~\citep{icse2026reproducibility}. The Moltbook incident (2026) exposed 1.5M API keys through vibe-coded services.

\paragraph{Latent execution.}
LPN (NeurIPS 2025)~\citep{macfarlane2025lpn}, COCONUT (Meta, 2024)~\citep{hao2024coconut}, LaSynth (NeurIPS 2021)~\citep{chen2021lasynth} demonstrate execution in representation space. ``Beyond Syntax'' (ICSE 2025)~\citep{beyondsyntax2025} shows LLMs develop language-independent code semantics. CWM (Meta, 2025)~\citep{metafair2025cwm} models code execution at token level.

\paragraph{Tokenization, naturalness, and dialects.}
Tokenizers introduce up to 15$\times$ cost unfairness across languages~\citep{petrov2024tokenizers}, with dialectal varieties suffering 10--15$\times$ tokenization premiums~\citep{tokenbiases2025}. Software is ``natural''---statistically predictable like NL~\citep{hindle2012naturalness}. Dialect-aware benchmarks~\citep{faisal2025swan} and studies of dialect bias in LLMs~\citep{bui2025dialect} reveal that the communicability gap extends within languages, not just across them.

\paragraph{Neuroscience and philosophy.}
Hyperalignment~\citep{haxby2011hyperalignment} and RSA~\citep{kriegeskorte2008rsa} show cross-brain representational convergence---the biological analogue of PRH---yet shared geometry does not resolve the private language problem~\citep{wittgenstein1953}. Mahowald et al.\ (2024)~\citep{mahowald2024dissociating} dissociate formal from functional linguistic competence in LLMs, paralleling our $\Zsem$/$\Zproc$ stratification. The modality gap in contrastive learning~\citep{liang2022modalitygap} persists even after alignment; closing it helps group-level but not instance-level tasks~\citep{grassucci2026modgap}, consistent with our prediction that $\Zsem$ converges while $\Zprag$ does not.

%% ============================================================
%% 8. CONCLUSION
%% ============================================================

\section{Conclusion}

PRH establishes that neural representations converge. We have argued that this is a beginning, not an end: the question that matters for the NL-code interface is not whether $Z$ exists, but whether it is \emph{communicable}.

Mathematics is often presented as evidence that communicability is achievable: Platonic ideals function as a cross-cultural medium despite divergent derivation paths and conventional notation. But this framing overstates mathematics' purity. Mathematics is itself a human-constructed system of meaning assignment---numerals carry cultural connotations that denotation cannot suppress (the avoidance of 4 in East Asian cultures, of 13 in Western ones, of 666 in Abrahamic traditions). Mathematical notation is culturally contingent (Arabic vs.\ Roman vs.\ CJK numerals; the late and uneven adoption of zero). What mathematics achieves is not denotation without human meaning, but a \emph{social convention to suppress cultural connotation at the point of formal reasoning}---a convention that NL, by its nature, cannot adopt. Mathematics achieves communicability not because it bypasses human interpretation, but because it disciplines it. Our pilot provides evidence that the NL-to-$Z$ encoding path does not inherit even this disciplined property: execution-level convergence (NL-code alignment, $R_{\text{code}} > 1$) coexists with description-level divergence (P2 failure). Representations can converge without the descriptions that access them converging---existence does not entail communicability.

This parallels challenges in neuroscience (representational similarity $\neq$ shared experience), philosophy (universals $\neq$ intersubjective access), and denotational semantics (shared denotation $\neq$ shared specification). Three practical consequences follow: disambiguation cost is conserved, the Chomsky wall retreats inward to verification, and $\Dtrain$ dependence requires explicit cultural parameterization. Mathematics achieved communicability for abstract reasoning---not by transcending human interpretation, but by socially disciplining it. Whether computation can achieve the same discipline for executable intent, and whether agent-to-agent $Z$-space communication can bypass the human interpretation bottleneck entirely, remain open questions with immediate architectural consequences for multi-agent system design.

%% ============================================================
%% LIMITATIONS (does not count toward page limit)
%% ============================================================

\section*{Limitations}

\textbf{PRH is a hypothesis, not a theorem.} PRH for code as a modality remains a conjecture. Our contribution is not to prove it but to derive consequences and provide initial tests. Recent evidence (language-independent code semantics~\citep{beyondsyntax2025}, cross-model transferability~\citep{chen2025crossmodel}) is consistent with convergence but not proof.

\textbf{Pilot measures description-level, not execution-level, convergence.} The P2 failure highlights this gap: NL embedding similarity is a proxy, not a direct test, of $\Zsem$ convergence. A stronger test would embed code alongside NL descriptions and measure alignment in code-trained representation spaces. We have begun this with the NL-code alignment experiment (UniXcoder), but a full-scale study across model families and scales is needed.

\textbf{The $Z$ stratification is a conceptual framework.} While the pilot provides initial evidence (P7 supported, P2 failure consistent with the description/execution distinction, P3 supporting cross-lingual separability), the stratification has not been validated with large-scale probing across multiple model families.

\textbf{Disambiguation bound is a design heuristic.} The information-theoretic formalization (Eq.~\ref{eq:conservation}) draws on Kolmogorov complexity, which is uncomputable in general. The bound principle guides system design but is not a formal theorem; it establishes a lower limit on specification cost, not a conservation law.

\textbf{Neuroscience connections are structural analogies.} We draw parallels between PRH convergence and Hyperalignment, and between the communicability gap and the private language problem. These are structural correspondences, not claims of mechanistic identity between artificial and biological neural networks.

\textbf{Representation-level execution is a research direction,} not a working system. The pieces exist (SONAR, LCM, LPN, constrained decoding) but have not been unified.

%% ============================================================
%% ETHICS STATEMENT (does not count toward page limit)
%% ============================================================

\section*{Ethics Statement}

This work identifies ethical risks inherent in representation-level computing systems that we believe the community should address proactively:

\textbf{Cultural bias in $Z$.} If shared representations are shaped by English-dominant training data, they may encode Anglo-American computational norms as ``neutral'' defaults (Appendix~\ref{app:ethics}). This constitutes a form of computational colonialism~\citep{phillipson1992} that affects non-English-speaking users disproportionately.

\textbf{Consent and opacity.} Users interacting with representation-level systems cannot inspect which cultural assumptions shaped the interpretation of their input. We advocate for cultural provenance tracking and explicit cultural parameterization (Appendix~\ref{app:ethics}, mitigations M1--M5).

\textbf{Feedback loops.} $Z$-projected outputs re-enter training distributions, potentially amplifying cultural biases through recursive training~\citep{shumailov2024collapse}.

Our pilot experiment uses publicly available multilingual sentence-transformer models and does not involve human subjects. All stimuli were created by the authors.

%% ============================================================
%% REFERENCES
%% ============================================================

\bibliography{references}
\bibliographystyle{colm2026_conference}

%% ============================================================
%% APPENDICES
%% ============================================================

\appendix

\section{Detailed Code Examples}\label{app:code}

\subsection{Wave 1: Keyword Substitution (Han --- Korean Rust)}

Han~\citep{han2025} maps Korean keywords to Rust syntax. The grammar remains Type-2; only terminal symbols change:

\begin{lstlisting}[language=Python,escapeinside={(*}{*)}]
func top_set(list: array<int>) -> array<int> {
    sort(list, reverse=true)[..3]
}
\end{lstlisting}

The parse tree is identical to its English-syntax equivalent. The programmer must still know Type-2 grammar rules---only the vocabulary changes.

\subsection{Wave 2: Structured NL (CoRE-style)}

CoRE~\citep{core2025} constrains NL into a parseable template format:

\begin{lstlisting}
TASK: Find three largest values
INPUT: list of numbers
STEPS:
  1. Sort the list in descending order
  2. Take the first three elements
OUTPUT: list of three numbers
\end{lstlisting}

This is a new Type-2 grammar that \emph{looks} like English. The keywords (TASK, INPUT, STEPS, OUTPUT) are non-terminals; the prose within each field is constrained to parseable templates. Making NL precise enough to execute requires removing the properties that make it natural.

\subsection{Wave 3: LLM-Mediated Translation}

Three runs of the same prompt ``Find the three largest values in the given list'' produce:

\begin{lstlisting}[language=Python]
# Run 1
sorted(lst, reverse=True)[:3]

# Run 2
import heapq
heapq.nlargest(3, lst)

# Run 3 (buggy)
[x for x in lst if x >= sorted(lst)[-3]]  # fails on duplicates
\end{lstlisting}

All three are valid interpretations; Run~3 contains a bug that manifests only with duplicate values. The non-determinism is not a bug in the LLM---it reflects genuine ambiguity in the NL input.

\subsection{The RAG Parallel}

Both RAG and NL$\to$Code map NL to structured output via learned representations. The critical difference: RAG tolerates approximate retrieval (cosine similarity 0.95 is acceptable), while code requires exact semantics (a single wrong token is a bug). This distinction maps directly to the convergence-determinism gap (\S3.7): $\delta > 0$ is acceptable for retrieval but not for execution.

\section{Training Distribution Dependence of $Z$}\label{app:dtrain}

\subsection{$Z$ as a Function of Training Data}

The shared latent structure is not determined by scale alone:
\[
Z = Z(\text{scale}, \Dtrain)
\]
Two models at the same scale trained on different corpora may converge to different $Z$. This dependence is not limited to culturally obvious domains---even ``purely computational'' operations carry implicit cultural defaults.

\subsection{Cultural Defaults in ``Purely Computational'' Domains}

\begin{itemize}[nosep]
  \item \textbf{Sort order conventions.} Default ascending in most Western languages; Japanese databases sometimes default to JIS code order; Chinese systems may default to stroke-count or pinyin order.
  \item \textbf{Null handling semantics.} SQL NULL $\neq$ Python None $\neq$ JavaScript undefined/null $\neq$ Rust Option::None. Each reflects a different cultural stance on ``absence.''
  \item \textbf{Number formatting.} \texttt{1,000.00} (US) vs.\ \texttt{1.000,00} (Germany) vs.\ \texttt{1,000.00} with \emph{man} (ten-thousand) grouping (Korea). A ``purely computational'' formatter encodes cultural convention.
  \item \textbf{Error semantics.} Python exceptions (ask forgiveness, not permission), Go explicit errors (check everything), C undefined behavior (trust the programmer). These are not just API differences---they encode cultural attitudes toward failure and responsibility.
\end{itemize}

\subsection{Training Distribution Divergence}

The clean $\Zsem$/$\Zproc$ separation is an idealization. In practice:
\[
\Zsem(\text{observed}) = \Zsem(\text{ideal}) + \mathrm{bias}(\Dtrain)
\]

The sovereign AI movement implicitly acknowledges $\Dtrain$ dependence:
\begin{itemize}[nosep]
  \item Korea's five-consortium initiative for Korean-native LLMs
  \item Japan's ABCI 3.0 for Japanese-native training
  \item SEA-LION for Southeast Asian languages
\end{itemize}
If $Z$ were truly Platonic---culture-invariant at all layers---culturally-native models would be unnecessary. The market is voting for Model~A.

\section{Tokenization Bottleneck and Spacing Analysis}\label{app:tokenization}

\subsection{Language Types and Tokenization Behavior}

\begin{table*}[h]
\centering
\caption{Tokenization behavior across language types.}
\label{tab:tokenization}
\small
\begin{tabularx}{\textwidth}{l l l L}
\toprule
\textbf{Language type} & \textbf{Spacing} & \textbf{Token boundary} & \textbf{Impact on $Z$} \\
\midrule
Space-delimited (English) & Mandatory & Aligned with words & Minimal---tokenizer matches linguistic units \\
Agglutinative (Korean) & Partial (\emph{ttieo-sseugi}) & Misaligned & Spacing variants produce different token sequences $\to$ different $Z$ \\
Unsegmented (Chinese) & None & Character/subword & Segmentation ambiguity leaks into $Z$ \\
Abjad (Arabic) & Space-delimited & Right-to-left + morphological & Diacritics affect tokenization \\
\bottomrule
\end{tabularx}
\end{table*}

\subsection{Korean Spacing Variants}

Korean spacing (\emph{ttieo-sseugi}) is notoriously inconsistent even among native speakers. Four variants of our running example:

\begin{enumerate}[nosep]
  \item \texttt{jueo-jin mogrog-eseo gajang keun gabs se gae-reul chaj-ara} (correct spacing)
  \item \texttt{jueo-jinmogrog-eseo gajangkeungabs segae-reul chaj-ara} (common informal)
  \item \texttt{jueo-jinmogrog-eseogajangkeungabssegae-reulchaj-ara} (no spaces)
  \item \texttt{ju eo-jin mog rog-eseo gajang keun gabs se gae-reul chaj-ara} (over-spaced)
\end{enumerate}

All four carry identical computational intent. A Korean speaker processes all four effortlessly; a tokenizer produces four different token sequences, potentially mapping to four different points in $Z$.

\subsection{Token Tax}

Petrov et al.\ (2024)~\citep{petrov2024tokenizers} document that non-Latin languages require 2--15$\times$ more tokens for the same semantic content. This ``token tax'' means: higher API cost for non-English users, reduced effective context window, and---most relevant to PRH---longer paths through the model, potentially affecting where in $Z$ the input lands.

\subsection{Prediction 7 Protocol}

P7 (Spacing Robustness) can be tested by measuring:
\[
\frac{d(Z(\text{variant}_i), Z(\text{variant}_j))}{d(Z(\text{sentence}_1), Z(\text{sentence}_2))} \ll 1
\]
where variants $i,j$ are spacing variants of the same sentence and sentences $1,2$ are semantically different. The ratio should be $> 10\times$ for a robust encoder.

\paragraph{The ideal encoder behaves like a Korean speaker.}
A Korean speaker maps all four spacing variants to the same meaning effortlessly. The ideal encoder should do the same---collapsing surface variation while preserving semantic distinction. Current tokenizers do not; character-level models (ByT5~\citep{xue2022byt5}, CANINE~\citep{clark2022canine}) partially address this but sacrifice other capabilities.

\section{Symbols, Emoji, and the Spectrum of Representational Directness}\label{app:emoji}

\subsection{The Representation Spectrum}

Between NL and $Z$ exists a spectrum of surface forms varying in cultural dependency and ambiguity:

\begin{table*}[h]
\centering
\caption{Spectrum of representational directness from NL to $Z$.}
\label{tab:spectrum}
\small
\begin{tabularx}{\textwidth}{l l l L}
\toprule
\textbf{Surface form} & \textbf{Cultural dependency} & \textbf{Ambiguity} & \textbf{Example} \\
\midrule
Natural language & High & High & ``find the largest value'' \\
Emoji/pictographs & Medium & Medium & Folded hands emoji: prayer / gratitude / high-five \\
Mathematical symbols & Low & Low & $\max(S)$, $|S| = 3$ \\
Code operators & Very low & Very low & \texttt{sorted(lst)[:3]} \\
$Z$ (continuous) & None (target) & None (target) & $z \in \mathbb{R}^d$ \\
\bottomrule
\end{tabularx}
\end{table*}

\subsection{Emoji as $\Zsem$ Experiments}

Emoji are natural experiments in cross-cultural semantic convergence:
\begin{itemize}[nosep]
  \item Near-universal recognition (the thumbs-up emoji means ``good'' in most contexts).
  \item Cultural variation that mirrors $Z$ stratification (the folded-hands emoji is prayer in Western contexts, gratitude in Japanese, high-five in some youth cultures).
  \item Composition parallels: emoji sequences compose meaning (face + heart = love), but compositionality is loose---paralleling NL rather than code.
\end{itemize}

Lu et al.\ (2016)~\citep{lu2016emoji} document usage patterns; Barbieri et al.\ (2016)~\citep{barbieri2016emoji} measure cross-cultural variation. The pattern matches $Z$ stratification: $\Zsem$ (core meaning) converges while $\Zprag$ (contextual usage) remains culturally specific.

\subsection{Mathematical Symbols: Pre-Neural $\Zsem$ Convergence}

Mathematical notation achieved $\Zsem$ convergence before neural networks existed. The symbol $\int$ means the same operation worldwide; the notation is conventional (Leibniz won over Newton) but the denotation is universal. This is direct evidence that Platonic convergence is possible---and that surface-form conventions persist by path dependency, not necessity.

\section{Ethical Implications of $\Dtrain$-Dependent $Z$}\label{app:ethics}

\subsection{Five Dimensions of Risk}

\begin{enumerate}[nosep]
  \item \textbf{Computational colonialism.} If $Z$ is trained predominantly on English data, English-derived computational norms become the ``neutral'' baseline. Non-English users must conform to Anglo-American defaults---sort order, error handling, null semantics---or accept degraded performance. This is Phillipson's (1992)~\citep{phillipson1992} linguistic imperialism extended to computational semantics.

  \item \textbf{The ``purely computational'' boundary.} The distinction between ``purely computational'' ($\Zsem$) and ``culturally inflected'' ($\Zproc$, $\Zprag$) is itself a cultural construction. What counts as ``purely computational'' depends on who draws the boundary. Risk assessment, prioritization, and scheduling all involve implicit value judgments that may appear ``computational'' from within one cultural framework.

  \item \textbf{Consent and transparency.} Users interacting with representation-level systems cannot see which cultural assumptions shaped the interpretation of their input. If ``find the largest value'' is interpreted using English-derived defaults for comparison and ordering, the Korean user has no visibility into this cultural translation.

  \item \textbf{Accountability gap.} When a representation-level system misinterprets culturally-mediated intent, who is responsible? The user (who expressed intent in NL)? The system (which mapped to $Z$)? The training data (which shaped $Z$)? The $\Dtrain$ dependence creates a ``responsibility gap''~\citep{matthias2004} where no single actor controls the cultural mapping.

  \item \textbf{Feedback loop risks.} $Z$-projected outputs re-enter the training distribution when users interact with and build upon system outputs. If $Z$ carries cultural bias, this bias is amplified through recursive training~\citep{shumailov2024collapse}, creating a feedback loop where dominant cultural norms become increasingly entrenched in $Z$.
\end{enumerate}

\subsection{Mitigation Paths}

\begin{description}[nosep]
  \item[M1: Cultural provenance tracking.] Annotate $Z$ projections with cultural assumptions: ``This interpretation assumes ascending sort order (Western convention).''
  \item[M2: Explicit cultural parameterization.] Make cultural defaults configurable: $Z(c)$ where $c$ specifies cultural context, rather than a single ``universal'' $Z$.
  \item[M3: Diverse $\Dtrain$ requirements.] Mandate training data cultural diversity metrics, analogous to demographic diversity requirements in dataset documentation~\citep{durmus2023opinions}.
  \item[M4: Cross-cultural validation.] Test representation-level systems across cultural contexts, not just languages---measuring whether the system's defaults align with each culture's expectations.
  \item[M5: User-inspectable cultural mapping.] Develop tools that allow users to see how their NL input was interpreted---which cultural defaults were applied, and how to override them.
\end{description}

\section{Pilot Experiment Design (Full Protocol)}\label{app:pilot}

\subsection{Stimuli}

\begin{itemize}[nosep]
  \item \textbf{Computational operations (50):} sorting, filtering, aggregation, string manipulation, mathematical operations, data transformation, search algorithms, set operations, type conversion, and formatting---operations with well-defined denotational semantics.
  \item \textbf{Judgment operations (50):} risk assessment, prioritization, categorization (subjective), recommendation, scheduling, resource allocation, quality assessment, sentiment-dependent filtering, fairness-constrained selection, and culturally-dependent formatting---operations where ``correct'' depends on cultural context.
  \item \textbf{Languages (5):} English, Korean, Chinese, Arabic, Spanish---chosen for typological diversity (SVO/SOV, space-delimited/agglutinative/unsegmented, Latin/Hangul/CJK/Arabic scripts).
  \item \textbf{Total stimuli:} $100 \times 5 = 500$ NL descriptions.
\end{itemize}

\subsection{Models}

\begin{itemize}[nosep]
  \item \textbf{SONAR} (Meta): Sentence-level embeddings, 200 languages, speech-capable.
  \item \textbf{StarCoder-2} / code-specific model: Code-trained embeddings for comparison.
  \item \textbf{Frontier models:} GPT-4, Claude, Gemini---extract hidden-state embeddings via API probing or use output-based similarity measures.
\end{itemize}

\subsection{Protocol}

\begin{enumerate}[nosep]
  \item \textbf{Embed} all 500 stimuli through each model.
  \item \textbf{Compute} pairwise cosine similarity matrices ($500 \times 500$ per model).
  \item \textbf{Calculate discriminability ratio} $R = d_{\text{inter}} / d_{\text{intra}}$ where:
    \begin{itemize}[nosep]
      \item $d_{\text{intra}}$: mean distance between same-operation descriptions across languages (e.g., ``sort'' in English vs.\ Korean).
      \item $d_{\text{inter}}$: mean distance between different-operation descriptions within the same language (e.g., ``sort'' vs.\ ``reverse'' in English).
    \end{itemize}
  \item \textbf{Compare} $R$ for computational vs.\ judgment operations.
\end{enumerate}

\subsection{Predictions}

\begin{itemize}[nosep]
  \item \textbf{P2:} $R > 1$ for computational operations (cross-lingual same-operation similarity $>$ within-language different-operation similarity).
  \item \textbf{P2 extension:} $R \approx 1$ or $R < 1$ for judgment operations.
  \item \textbf{P1:} $R$ increases with model scale (test across model sizes if available).
  \item \textbf{P6:} $R$ varies across models with different $\Dtrain$ (GPT-4 vs.\ HyperCLOVA X vs.\ Qwen).
\end{itemize}

\subsection{Controls}

\begin{itemize}[nosep]
  \item Random baseline: $R \approx 1$ for randomly paired operations.
  \item Paraphrase control: same-language paraphrases should have $R \gg 1$.
  \item Translation quality: use professional translations, not MT, to isolate representation effects from translation artifacts.
\end{itemize}

\subsection{Extensions}

\paragraph{P6: $\Dtrain$ sensitivity.}
Compare embeddings from models with known $\Dtrain$ differences:
\begin{itemize}[nosep]
  \item GPT-4 (English-dominant $\Dtrain$)
  \item HyperCLOVA X (Korean-augmented $\Dtrain$)
  \item Qwen (Chinese-augmented $\Dtrain$)
\end{itemize}
Measure whether $\Zsem$ agreement correlates with $\Dtrain$ overlap.

\paragraph{P7: Spacing robustness.}
Test Korean spacing variants (4 per sentence, see Appendix~\ref{app:tokenization}):
\begin{itemize}[nosep]
  \item Correct spacing
  \item Common informal spacing
  \item No spaces
  \item Over-spaced
\end{itemize}
Measure $d(\text{spacing variants}) / d(\text{semantic variants})$; target: $> 10\times$.

\paragraph{P2 extension: Dialect continuum.}
Test whether the communicability gap is continuous rather than categorical by computing $R$ across dialectal variants within each language:
\begin{itemize}[nosep]
  \item \textbf{Korean:} Standard (Seoul) / Gyeongsang / Jeolla / Jeju
  \item \textbf{Arabic:} MSA / Egyptian / Gulf / Levantine
  \item \textbf{Chinese:} Mandarin / Cantonese / Min Nan
  \item \textbf{English:} American / British / Indian / AAVE
  \item \textbf{Spanish:} Castilian / Mexican / Rioplatense
\end{itemize}
Prediction: $R_{\text{within-dialect}} > R_{\text{cross-dialect}} > R_{\text{cross-lingual}}$, forming a continuum. Jeju Korean and Levantine Arabic---often classified as borderline separate languages---should approach cross-lingual $R$ values, exposing the arbitrariness of discrete language boundaries in embedding space. This directly tests whether modality incompleteness scales with linguistic distance.

\paragraph{P8: Keyboard transliteration fragility.}
Convert all Korean stimuli to their English-keyboard transliterations (e.g., ``정렬하라'' $\to$ ``wjddufgkfk'') using the standard KS X 5002 Dubeolsik layout mapping. Compute $d(\text{transliterated}, \text{correct})$ and compare against $d(\text{spacing variant}, \text{correct})$ from P7. Prediction: transliteration drift $\gg$ spacing drift $\gg$ punctuation drift, establishing a hierarchy of surface-variant destructiveness. This also applies to other keyboard layouts: Arabic (ASMO 449), Chinese Pinyin input without tone marks, Russian (JCUKEN). If any model achieves low transliteration drift, it has learned keyboard-layout invariance---an implicit signal that the model encodes intent rather than surface form.

\paragraph{Voice input extension.}
Test speech input via SONAR's speech encoder---spacing is absent in speech, so the spacing problem should disappear entirely. Compare $R$ for text vs.\ speech input.

\subsection{Feasibility}

\begin{itemize}[nosep]
  \item \textbf{Estimated cost:} $<$\$1{,}500 (API calls for frontier models; SONAR is open-source).
  \item \textbf{Timeline:} 3--4 weeks (stimulus preparation: 1 week; embedding + analysis: 1 week; writing: 1--2 weeks).
  \item \textbf{Compute:} Single GPU for SONAR; API access for frontier models.
  \item \textbf{Data:} Stimuli created by bilingual researchers; professional translations for non-native languages.
\end{itemize}

\subsection{Full Stimuli List}

\paragraph{Computational operations (50).}
\small
\begin{enumerate}[nosep]
\item sort ascending \item find max \item filter positive \item reverse list \item count elements \item sum \item deduplicate \item top-3 largest \item mean \item sort descending \item concatenate strings \item uppercase \item split by spaces \item replace characters \item string length \item absolute value \item cube \item modulo 7 \item square root \item GCD \item set union \item set intersection \item set difference \item extract keys \item merge dicts \item transpose matrix \item flatten list \item zip lists \item find index \item contains check \item all positive \item any negative \item int to string \item round to 2 decimals \item find min \item median \item slice [2:5] \item filter even \item frequency count \item double all \item tree depth \item is palindrome \item to binary \item cumulative sum \item product \item range 1--10 \item character count \item is prime \item sort by length \item matrix multiply
\end{enumerate}

\paragraph{Judgment operations (50).}
\begin{enumerate}[nosep]
\item prioritize by importance \item evaluate quality \item summarize document \item recommend best option \item assess risk \item categorize into groups \item select key customers \item rank by relevance \item flag problematic entries \item adjust tone \item determine sentiment \item decide appropriateness \item classify by urgency \item assess credibility \item simplify explanation \item compare proposals \item suggest compromise \item evaluate ethicality \item generate creative title \item assign to team members \item allocate budget fairly \item select best candidate \item write constructive feedback \item rate trustworthiness \item determine relevant context \item identify harmful content \item predict success likelihood \item translate preserving nuance \item scope first release \item draft apology \item detect bias \item write motivational message \item suggest fair price \item suggest UI improvements \item resolve requirement conflict \item write persuasive argument \item respond with empathy \item develop marketing strategy \item diagnose root cause \item decide on second chance \item handle sensitive topic \item propose innovative solution \item determine acceptable error rate \item craft data narrative \item rewrite rejection politely \item choose feature name \item analyze speed/accuracy tradeoff \item set trust level \item determine launch timing \item maintain or retire legacy system
\end{enumerate}
\normalsize

Each operation is described in English, Korean, Chinese, Arabic, and Spanish by native-speaker computational linguists ($100 \times 5 = 500$ base stimuli). Dialectal variants add up to 4 variants per language (${\sim}1{,}800$ total stimuli including dialect extensions). Full multilingual descriptions are available in the experiment repository.

%% ============================================================
%% REPRODUCIBILITY STATEMENT (COLM requirement, does not count toward page limit)
%% ============================================================

\section*{Reproducibility Statement}

All embedding experiments use publicly available pre-trained models (MiniLM-L12, E5-small/base/large, BGE-M3) in inference mode with default settings. The embedding computation is deterministic (no sampling or temperature parameters). Statistical significance tests use 10,000-permutation tests. Probing classifier results use 5-fold cross-validation with mean accuracy reported. All stimuli (100 operations $\times$ 5 languages = 500 base descriptions, plus dialectal variants totaling approximately 1,800 stimuli) were created by the authors and are available in the experiment repository. Total compute time: approximately 4 hours on macOS with Apple Silicon (M-series), CPU inference only.

\end{document}
